\documentclass[11pt, a4paper]{article}
\usepackage[a4paper, margin=2.5cm]{geometry}
\usepackage{fontspec}
\usepackage[english, russian, bidi=basic, provide=*]{babel}
\babelfont{rm}{Noto Sans}
\usepackage{amsmath}
\usepackage{booktabs}
\usepackage{hyperref}

\title{\textbf{Отчет по лабораторной работе №2}\\Документирование и рефакторинг кода}
\author{[Ваше ФИО / Номер ИСУ]}
\date{\today}

\begin{document}

\maketitle

\section{Цель работы}
Освоение практических навыков рефакторинга плохо документированного кода, внедрения стандартизированных блоков документации (docstrings), ведения проектной документации в формате Markdown в каталоге \texttt{docs/}, а также управления версиями с использованием Git (ветвление, коммиты, слияние веток).

\section{Описание выполненных задач}

В рамках лабораторной работы был исследован файл \texttt{mysterious\_file.py}, содержащий функции с неинтуитивными именами и отсутствующей или фрагментарной документацией. Был выполнен полный рефакторинг:
\begin{itemize}
    \item Переименованы функции и переменные для обеспечения читаемости и прозрачности бизнес-логики.
    \item Добавлены типы аргументов и возвращаемых значений (Type Hinting).
    \item Для каждой функции составлены подробные строки документации (docstrings) в соответствии со стандартом, включающие описание параметров (\texttt{Args}), возвращаемых значений (\texttt{Returns}) и примеров использования.
\end{itemize}

\section{Перечень рефакторенных функций}

\begin{itemize}
    \item \textbf{\texttt{make\_tea}}: Приготовление чая (вода и сахар).
    \item \textbf{\texttt{calculate\_change}}: Расчет сдачи при покупке или проверка достаточности средств.
    \item \textbf{\texttt{is\_weekend}}: Проверка, является ли переданный день недели выходным.
    \item \textbf{\texttt{is\_even}}: Проверка четности числа.
    \item \textbf{\texttt{filter\_numbers\_divisible\_by\_three}}: Фильтрация чисел в диапазоне, кратных трем.
    \item \textbf{\texttt{calculate\_sentiment\_score}}: Оценка тональности текста с учетом знаков препинания.
    \item \textbf{\texttt{decimal\_to\_binary\_string}}: Конвертация десятичного числа в двоичную строку.
\end{itemize}

\section{Структура проектной документации}

В корне проекта был создан каталог \texttt{docs/}, содержащий файл \texttt{documentation.md}. В нем зафиксированы:
\begin{enumerate}
    \item Подробное описание каждой функции с указанием сигнатур и примеров вызова.
    \item История изменений репозитория и ветки разработки.
\end{enumerate}

\section{История коммитов и Git-воркфлоу}

Разработка велась с соблюдением требований Git:
\begin{enumerate}
    \item Создана вспомогательная ветка \texttt{doc\_<ИСУ>} для документирования.
    \item Выполнены осмысленные коммиты с фиксацией промежуточных результатов.
    \item Создана персональная ветка разработки \texttt{main\_features\_<ИСУ>}, в которую выполнены слияние (\texttt{git merge}) и последующий Pull Request на GitHub.
\end{enumerate}

\section{Заключение}
В ходе выполнения лабораторной работы были закреплены навыки командной разработки, документирования кода на Python и оформления технической отчетности с использованием систем верстки (LaTeX).

\end{document}